\section{Imperfect advice, joint graph evidence and retention guarantees}
\label{exec17:advice-development}

A banking analyst may know that a diagram remains ambiguous while still
having useful information about the disputed arrows. An account officer can
describe which event usually precedes another; a procedure can explain a
permitted action; a language model can suggest which interpretation to
investigate. The difficult question is how much graph uncertainty this advice
can remove without concealing its mistakes. Requiring an adviser to be
infallible discards potentially useful information. Treating each confident
answer as a causal fact discards potentially correct graphs. A proposal needs
an explicit rule between these extremes, together with a statement of the
probability that the rule actually controls.

Long and colleagues formulate this problem using language models as imperfect
experts. Their 2023 first version starts from the true Markov equivalence class
and uses metadata-based judgments to orient its ambiguous edges. A prior over
legal graphs and an adviser-error likelihood produce joint posterior
probabilities; greedy choices trade a smaller class against the modeled risk
of excluding its true member. This is a useful methodological alternative to
hard acceptance of an elicited diagram. The known starting class and the
specified adviser model are substantive premises, however, rather than
consequences of a natural-language interface.\citep{exec17_imperfect_experts}

\subsection{The starting class and the experiment that supplies advice}

Let $G$ be a directed acyclic graph, and let $\mathcal M_0$ be a finite class of
graphs with the same skeleton and unshielded colliders. For an ambiguous pair
$j$, write $O_j(G)\in\{0,1\}$ for its orientation in $G$ and $E_j$ for the
adviser's proposed orientation. These binary judgments concern an existing
adjacency; they do not decide whether the two variables are the right
measurements, whether an edge is absent, or whether a hidden common cause
should replace it. In particular, the method's pair question does not provide
a general test of the bank's variable ontology or latent-variable assumptions.

Under a uniform prior on $\mathcal M_0$ and a conditionally independent adviser
with constant error $\epsilon$, the likelihood and posterior are
\begin{equation}
\begin{split}
 L(g;e)&=\prod_{j=1}^{u}(1-\epsilon)^{\mathbf 1\{O_j(g)=e_j\}}
                   \epsilon^{\mathbf 1\{O_j(g)\ne e_j\}},\\
 \pi(g\mid e)&=\frac{\mathbf 1\{g\in\mathcal M_0\}L(g;e)}
                        {\sum_{h\in\mathcal M_0}L(h;e)}.
\end{split}
\label{exec17:advice-posterior}
\end{equation}
Here $u$ is the number of ambiguous pairs, $e$ is the observed advice vector,
and $\mathbf 1$ denotes an indicator. The likelihood factorization is a claim
about how adviser errors arise conditional on the graph. The posterior does
not factorize over edges, because the legal graph class couples their
orientations. Long et al. make this distinction explicitly. A cycle or a new
unshielded collider can make a collection of individually plausible answers
jointly impossible.\citep{exec17_imperfect_experts}

A retained set $\mathcal R(e)\subseteq\mathcal M_0$ can be required to satisfy
\begin{equation}
 \sum_{g\in\mathcal R(e)}\pi(g\mid e)\geq 1-\eta,
 \qquad 0\leq\eta\leq1.
\label{exec17:advice-retention}
\end{equation}
The tolerance $\eta$ describes excluded posterior mass under this model.
The paper considers greedy choices based on the resulting class size or
retention risk, propagating consequences of each chosen orientation. It does
not establish that its greedy rule solves every global minimum-cardinality
retention problem. More fundamentally, the word ``probability'' in
\eqref{exec17:advice-retention} needs its probability model. A posterior over
graphs is useful conditional uncertainty; it is not automatically a
confidence statement for every fixed bank graph and every data-dependent
starting class.

\subsection{A small legal class exposes dependence between answers}

Consider the chain skeleton $A-B-C$, with no unshielded collider. Its three
legal DAGs are $g_1:A\to B\to C$, $g_2:A\leftarrow B\to C$, and
$g_3:A\leftarrow B\leftarrow C$. Code a forward edge as one, so their
orientation vectors are $(1,1),(0,1),(0,0)$. The fourth vector $(1,0)$ would
create $A\to B\leftarrow C$ and is absent from this equivalence class.
For $\epsilon=1/10$, advice $(1,1)$ gives
\begin{equation}
 \bigl(\pi(g_1\mid e),\pi(g_2\mid e),\pi(g_3\mid e)\bigr)
       =\frac{1}{91}(81,9,1).
\label{exec17:advice-chain}
\end{equation}
This calculation follows by normalizing the likelihoods
$(81,9,1)/100$. It is an exact project example, independently enumerated over
the legal graphs, rather than a reported result of the paper's software.
\citep{exec17_advice_validity_exact}

The posterior mass of $A\to B$ is $81/91$, approximately $0.8901$, while that
of $B\to C$ is $90/91$, approximately $0.9890$. At a posterior retention
threshold of $0.9$, the latter orientation can be retained by keeping
$\{g_1,g_2\}$; keeping only $g_1$ fails the threshold. Two answers generated by
an adviser with marginal error $0.1$ have not produced a graph with posterior
probability $0.9$. The legal-class restriction and joint evidence determine
the result. If instead the advice is the incompatible vector $(1,0)$, the
posterior becomes $(9,1,9)/19$. Each advised orientation then has posterior
mass $9/19$. Disagreement with the graph constraints provides a reason to
reduce confidence, not a license to force the fourth, inadmissible graph into
the class.

For a bank, this example illustrates why one cannot independently approve an
``advice causes activity'' arrow and a ``credit draw causes activity'' arrow
and assume that both fit the diagram already justified by data. The variables
need precise definitions and clocks before that particular class is even
plausible. If documented time ordering rules out one orientation, it is
additional structural information and should change the class explicitly.
It should not be disguised as stronger language-model confidence. Conversely,
if a credible account requires a hidden common cause, the three-DAG class is
the wrong model family; redistributing its posterior cannot repair the missing
graph.

\subsection{Repeated prompts and the missing adviser law}

The constant-error model is attractive because it makes the calculation
transparent. Its applicability to a language model is an empirical question.
Long et al. score the two directional answer tokens and normalize their log
probabilities with a softmax. Their fixed prompt gives deterministic output;
the experiments randomize among causal verbs to introduce variation. This
historical procedure concerns text-davinci-002 and text-davinci-003. A softmax
score is a model preference over the presented alternatives, while the needed
likelihood concerns error relative to a true orientation. Agreement between
the two quantities requires evidence.\citep{exec17_imperfect_experts}

The distinction is especially visible when an assistant asks the same question
repeatedly. Let the true direction $G$ have equal prior probabilities and let
one latent answer $S$ agree with $G$ with probability $0.9$. Suppose five calls
all repeat $S$. The correct joint model gives
\begin{equation}
 \Pr(G=1\mid E_1=\cdots=E_5=1)=0.9,
 \qquad
 \widetilde\Pr(G=1\mid E_1=\cdots=E_5=1)
        =\frac{9^5}{1+9^5}=\frac{59049}{59050}
\label{exec17:advice-repetition}
\end{equation}
for, respectively, the repeated-signal law and an incorrectly factorized
five-signal law. Enumerating the graph and latent-answer joint distribution
reproduces both calculations. The second number is approximately $0.999983$;
no corresponding new information entered the first experiment.
\citep{exec17_advice_validity_exact}

Changing a verb, sampling a different completion or asking a second model does
not by itself establish conditional independence. Models can share training
material, prompts can share a misconception, and several retrieved passages
can repeat one source. A practical adviser model may therefore need grouped
errors, variable-pair difficulty, abstention, asymmetric mistakes and dependence
on the wording or supplied documents. These are proposed modeling extensions,
not properties already verified for the bank. An alternative is to keep
elicited answers as soft sensitivity inputs and report how conclusions change
as their influence varies. Another is to use the answers only to prioritize
human investigation while leaving the admissible graph class intact.

\subsection{Posterior retention and coverage of a learned class}

Suppose the starting class is instead learned from observations $D$, yielding
$\widehat{\mathcal M}(D)$. If the true graph $g_\star$ is absent from that
class, every refinement that returns a subset also excludes it. Even a
perfectly computed adviser posterior cannot recover an omitted member. A
frequentist composition would need an additional argument. For example, if
\begin{equation}
\begin{split}
 \Pr_{g_\star}\{g_\star\notin\widehat{\mathcal M}(D)\}&\leq\alpha,\\
 \Pr_{g_\star}\{g_\star\notin\mathcal R(D,E)
       \mid g_\star\in\widehat{\mathcal M}(D)\}&\leq\eta,
\end{split}
\label{exec17:advice-coverage-premises}
\end{equation}
then the law of total probability gives
\begin{equation}
 \Pr_{g_\star}\{g_\star\notin\mathcal R(D,E)\}
       \leq\alpha+(1-\alpha)\eta\leq\alpha+\eta.
\label{exec17:advice-coverage-composition}
\end{equation}
The second premise concerns the actual joint data/advice experiment, including
selection of the class. It cannot be replaced by a pointwise posterior-mass
calculation. The displayed bound is a project derivation under the two stated
premises, not a new qualification result for the published method or the
proposed library.

This distinction affects downstream monetary decisions. If a policy is chosen
after selecting a graph, uncertainty about the retained graph and uncertainty
about its causal parameters both matter. For a policy $a$, a proposed robust
decision could examine the smallest identified value across retained graphs
and parameter laws. Such a value set must retain any nonidentification within
a graph; a posterior orientation does not supply missing intervention data.
A posterior expected-value decision is a different alternative, conditional
on the chosen graph and parameter priors. The proposal should present those
two decision principles with their assumptions and compare the actions they
recommend, rather than silently turning either into a bank guarantee.

\subsection{What to reuse and what a prototype must still establish}

The published experiments provide a useful warning as well as a useful method.
They derive the starting MEC from each benchmark's true graph, repeat the
synthetic-error experiments five times and the language-model experiments
eight times, and report failures associated with ALARM metadata ambiguity and
text-davinci-003 overconfidence. These findings motivate sensitivity to the
adviser law. They do not evaluate discovery errors in a learned starting class
or establish performance for a current banking ontology and model. The
paper's code and codebook links identify candidates for further inspection;
the paper alone does not establish their license, current dependencies or
runtime behavior.\citep{exec17_imperfect_experts}

The proposal can reuse the conceptual separation between a data-supported
class and fallible directional advice, together with joint consistency and
retention calculations. It should reject a default that hard-deletes every
graph contradicted by an uncalibrated answer. Direct code adoption or a port
needs a pinned source and rights review, agreement with independent small-class
references, and a declared regime in which enumeration or an approximation is
feasible. A scalable approximate posterior also needs an error statement;
having a graph-search API does not settle its approximation quality.

A public prototype can test these issues without expensive bank experiments.
It can start from stipulated graph classes, preserve independent graph truth,
separate correct-class and misspecified-class cases, vary adviser dependence,
and compare hard constraints, soft priors, abstention and unchanged-class
baselines. Renamed variables and genuinely unfamiliar mechanisms are useful
additional cases, but they test different generalizations. Report graph
retention, class size and downstream action loss separately. Such experiments
would establish whether a particular implementation deserves a banking trial;
the exact examples above establish why those comparisons are necessary.

\label{exec17:advice-development-end}
